\documentclass[11pt]{nyjm}
\numberwithin{equation}{section}

\title{Certified Bounds for Prime-Indexed Operators and a Trace Identity}
\author{Ryan O. Van Gelder}
\address{Citizen Gardens, 189 Private Drive 123, Crown City, Ohio, 45623}
\email{info@citizengardens.org}
\thanks{Thanks to pretty much everyone!}
\date{October 31, 2025}

\begin{document}



\begin{abstract}
We establish unconditional boundedness, Hilbert--Schmidt, and trace-class properties for a class of prime-indexed operators $L_\Lambda$ with explicit constants. For the Laplace kernel case with parameter $\theta>1$, the operator is positive trace-class, enabling a rigorous trace identity for analytic functions with explicit remainder bounds.
\end{abstract}
\maketitle
\section{Setup}
Let $\Pset$ be the primes. Fix $\theta\ge 1$ and $c>0$. Define $L_\Lambda:\ell^2(\Pset)\to\ell^2(\Pset)$ by
\[
(L_\Lambda x)_p=\sum_{q\in \Pset}\frac{e^{-c|\log p-\log q|}}{(pq)^{\theta/2}}\,x_q.
\]
Write $\Sone$ and $\Stwo$ for trace and Hilbert--Schmidt classes. Let $D=\mathrm{diag}(p^{-\theta})$ and $R=L_\Lambda-D$.

\section{Certified Theorems}
\begin{theorem}[T1: Boundedness]
\label{thm:T1}
If $\theta\ge 1$ and $c>|1-\theta/2|$, then
\[
\|L_\Lambda\| \le \frac{1}{\log 2}\cdot \frac{2c}{c^2-(1-\theta/2)^2}.
\]
Note that $\|D\|=2^{-\theta}$.
\end{theorem}

\begin{theorem}[T2: Hilbert--Schmidt]
\label{thm:T2}
If $\theta>1$, then
\[
\|L_\Lambda\|_{\Stwo}^2 \le \frac{2^{\,2(1-\theta)}}{(\theta-1)(2c+\theta-1)(\log 2)^2}.
\]
\end{theorem}

\begin{lemma}[Positivity]
\label{lem:PSD}
The function $\kappa(t)=e^{-c|t|}$ is positive definite on $\R$ by Bochner's theorem \cite[Theorem IX.11.13]{DunfordSchwartz}, witnessed by 
$ e^{-c|t|} = \int_{\mathbb{R}} e^{it\xi} \mu(d\xi) $ 
with $ d\mu(\xi) = \frac{c}{\pi(c^2+\xi^2)} d\xi $. 
Hence the matrix $\big(\kappa(\log p-\log q)\big)_{p,q}$ is PSD. 
Since $\big(p^{-\theta/2}q^{-\theta/2}\big)_{p,q}$ is rank-one PSD, 
the kernel $K$ is PSD by the Schur product theorem \cite[Theorem 5.2.1]{HornJohnson}.
\end{lemma}

\begin{theorem}[T3$'$: Trace Class for the Laplace kernel]
\label{thm:T3prime}
If $\theta>1$, then $L_\Lambda\ge0$ and $L_\Lambda\in \Sone$ with
\[
\|L_\Lambda\|_{\Sone}=\mathrm{Tr}(L_\Lambda)=\sum_{p\in \Pset}\frac{1}{p^{\theta}}<\infty.
\]
\end{theorem}
\begin{proof}
By Lemma~\ref{lem:PSD}, $L_\Lambda\ge0$. For positive operators, $T\in \Sone$ if and only if $\sum_n\langle Te_n,e_n\rangle<\infty$ for some orthonormal basis. With the standard basis, $\langle L_\Lambda e_p,e_p\rangle=K(p,p)=p^{-\theta}$. Since $\sum_{p}p^{-\theta}\le \sum_{n\ge2}n^{-\theta}<\infty$ for $\theta>1$, the claim follows and $\|L_\Lambda\|_{\Sone}=\mathrm{Tr}(L_\Lambda)$.
\end{proof}

\section{Trace Identity and a Stable Remainder Bound}
Let $\phi$ be analytic on a domain $\Omega\supset \sigma(L_\Lambda)\cup \sigma(D)$. Then $\phi(L_\Lambda)$ and $\phi(D)$ are defined by holomorphic functional calculus.

\begin{theorem}[Trace identity with resolvent bound]
\label{thm:trace}
For analytic $\phi$ on $\Omega\supset \sigma(L_\Lambda)\cup \sigma(D)$,
\[
\mathrm{Tr}\,\phi(L_\Lambda)=\sum_{p\in \Pset}\phi(p^{-\theta})+\mathrm{Tr}\big(\phi(L_\Lambda)-\phi(D)\big),
\]
and $ \phi(L_\Lambda)-\phi(D)\in \Sone$.
Moreover, for any smooth Jordan curve $\Gamma\subset\Omega$ enclosing $\sigma(L_\Lambda)\cup \sigma(D)$,
\[
\|\phi(L_\Lambda)-\phi(D)\|_{1}\le \frac{1}{2\pi}\,\operatorname{length}(\Gamma)\,
\sup_{z\in\Gamma}|\phi(z)|\,
\sup_{z\in\Gamma}\|(zI-L_\Lambda)^{-1}\|\,
\sup_{z\in\Gamma}\|(zI-D)^{-1}\|\,
\|R\|_{1}.
\]
In particular, with the circle $|z|=r>\max\{\|L_\Lambda\|,2^{-\theta}\}$,
\[
\|\phi(L_\Lambda)-\phi(D)\|_{1}\le
\frac{r}{(r-\|L_\Lambda\|)(r-2^{-\theta})}\,
\sup_{|z|=r}|\phi(z)|\;\|R\|_{1}.
\]
\end{theorem}
\begin{proof}
By Cauchy's formula and the resolvent identity,
\[
\phi(L_\Lambda)-\phi(D)=\frac{1}{2\pi i}\int_\Gamma \phi(z)\,(zI-L_\Lambda)^{-1}(L_\Lambda-D)(zI-D)^{-1}\,dz,
\]
which is trace-class since $R\in \Sone$. Take trace and norms to obtain the bounds.
\end{proof}

\section{Remarks and Constraints}
\begin{itemize}
\item If $\kappa$ is replaced by a non--positive-definite kernel, use factorization $L_\Lambda=AB$ with $A,B\in\Stwo$ to get $L_\Lambda\in \Sone$ under an exponential moment for $\kappa$ and $\theta>1$.
\item No claims on eigenvalue pairing with $\{p^{-\theta}\}$.
\end{itemize}

\begin{thebibliography}

\bibitem[DS63]{DS63} Dunford, N.; Schwartz, J.T. \textit{Linear Operators. Part II: Spectral Theory.} \mrev{MR0188745 (32 \#6184)} \zbl{0128.09803} Wiley, New York, 1963.
 
\bibitem[HJ13]{HJ13} Horn, R.A.; Johnson, C.R. \textit{Matrix Analysis,} 2nd ed. \mrev{MR3091739 (2013h:15001)} \zbl{1267.15001} Cambridge Univ. Press, Cambridge, 2013.
 
\bibitem[Sim05]{Sim05} Simon, B. \textit{Trace Ideals and Their Applications,} 2nd ed. \mrev{MR2154153 (2006f:47048)} \zbl{1074.47001} \doi{10.1090/surv/120} American Mathematical Society, Providence, RI, 2005.
\end{thebibliography}

\end{document}